% Môn: Vật lý
% Lớp: 11
% Chương: Chương I. Dao động
% Bài: L11C1 Bài 4. Bài tập về dao động điều hòa
% Số câu: 75
% App đọc file này trực tiếp — sửa rồi Commit trên GitHub, bấm Đồng bộ GitHub .tex

% L11C1 Bài 4. Bài tập về dao động điều hòa



% ===== Câu 13 =====
% ID: L11C1B4-04-DS
% Mức: TH
\dangbt{Viết phương trình dao động điều hòa.}
\begin{ex}
% ID: L11C1B4-04-DS
% Mức: TH
Vật có biên độ $4\,\text{cm}$, chu kì $1\,\text{s}$ và tại $t=0$ qua vị trí cân bằng theo chiều dương. Xét phương trình $x=4\cos(2\pi t/1-\pi/2)$ cm.
\choiceTF
{\True Biên độ là $4\,\text{cm}$.}
{\True Chu kì là $1\,\text{s}$.}
{\True Tại $t=0$, li độ bằng 0.}
{\True Tại $t=0$, vận tốc dương.}
\end{ex}

% ===== Câu 14 =====
% ID: L11C1B4-04-TL
% Mức: TH
\begin{bt}
% ID: L11C1B4-04-TL
% Mức: TH
Vật dao động điều hòa có biên độ $4\,\text{cm}$, chu kì $1\,\text{s}$; lúc $t=0$ đi qua vị trí cân bằng theo chiều âm. Viết phương trình li độ dạng cos.
\end{bt}

% ===== Câu 17 =====
% ID: L11C1B4-05-DS
% Mức: VD
\begin{ex}
% ID: L11C1B4-05-DS
% Mức: VD
Vật có biên độ $6\,\text{cm}$, chu kì $2\,\text{s}$ và tại $t=0$ qua vị trí cân bằng theo chiều dương. Xét phương trình $x=6\cos(2\pi t/2-\pi/2)$ cm.
\choiceTF
{\True Biên độ là $6\,\text{cm}$.}
{\True Chu kì là $2\,\text{s}$.}
{\True Tại $t=0$, li độ bằng 0.}
{\True Tại $t=0$, vận tốc dương.}
\end{ex}

% ===== Câu 18 =====
% ID: L11C1B4-05-TL
% Mức: VD
\begin{bt}
% ID: L11C1B4-05-TL
% Mức: VD
Vật dao động điều hòa có biên độ $6\,\text{cm}$, chu kì $2\,\text{s}$; lúc $t=0$ đi qua vị trí cân bằng theo chiều âm. Viết phương trình li độ dạng cos.
\end{bt}

% ===== Câu 21 =====
% ID: L11C1B4-06-DS
% Mức: VD
\begin{ex}
% ID: L11C1B4-06-DS
% Mức: VD
Vật có biên độ $8\,\text{cm}$, chu kì $4\,\text{s}$ và tại $t=0$ qua vị trí cân bằng theo chiều dương. Xét phương trình $x=8\cos(2\pi t/4-\pi/2)$ cm.
\choiceTF
{\True Biên độ là $8\,\text{cm}$.}
{\True Chu kì là $4\,\text{s}$.}
{\True Tại $t=0$, li độ bằng 0.}
{\True Tại $t=0$, vận tốc dương.}
\end{ex}

% ===== Câu 22 =====
% ID: L11C1B4-06-TL
% Mức: VD
\begin{bt}
% ID: L11C1B4-06-TL
% Mức: VD
Vật dao động điều hòa có biên độ $8\,\text{cm}$, chu kì $4\,\text{s}$; lúc $t=0$ đi qua vị trí cân bằng theo chiều âm. Viết phương trình li độ dạng cos.
\end{bt}

% ===== Câu 25 =====
% ID: L11C1B4-07-DS
% Mức: TH
\dangbt{Xác định các đại lượng dựa vào đồ thị dao động ($x-t$, $v-t$, $a-t$).}
\begin{ex}
% ID: L11C1B4-07-DS
% Mức: TH
Đồ thị $x-t$ cho biên độ $8\,\text{cm}$, chu kì $1\,\text{s}$ và ban đầu ở biên dương.
\choiceTF
{\True Biên độ $8\,\text{cm}$.}
{\True Tần số $1\,\text{Hz}$.}
{\True Vận tốc ban đầu bằng 0.}
{\True Có thể viết $x=8\cos(2\pi t/1)$ cm.}
\end{ex}

% ===== Câu 26 =====
% ID: L11C1B4-07-TL
% Mức: TH
\begin{ex}
% ID: L11C1B4-07-TL
% Mức: TH
Đồ thị li độ–thời gian cho biết $A=4$ cm, hai đỉnh dương liên tiếp cách nhau $1\,\text{s}$ và tại $t=0$ vật ở biên dương. Xác định $T,f,\omega$ và viết phương trình dao động.
\choiceTF

\end{ex}

% ===== Câu 29 =====
% ID: L11C1B4-08-DS
% Mức: VD
\begin{ex}
% ID: L11C1B4-08-DS
% Mức: VD
Đồ thị $x-t$ cho biên độ $9\,\text{cm}$, chu kì $2\,\text{s}$ và ban đầu ở biên dương.
\choiceTF
{\True Biên độ $9\,\text{cm}$.}
{\True Tần số $0.5\,\text{Hz}$.}
{\True Vận tốc ban đầu bằng 0.}
{\True Có thể viết $x=9\cos(2\pi t/2)$ cm.}
\end{ex}

% ===== Câu 30 =====
% ID: L11C1B4-08-TL
% Mức: VD
\begin{ex}
% ID: L11C1B4-08-TL
% Mức: VD
Đồ thị li độ–thời gian cho biết $A=6$ cm, hai đỉnh dương liên tiếp cách nhau $2\,\text{s}$ và tại $t=0$ vật ở biên dương. Xác định $T,f,\omega$ và viết phương trình dao động.
\choiceTF

\end{ex}

% ===== Câu 33 =====
% ID: L11C1B4-09-DS
% Mức: VD
\begin{ex}
% ID: L11C1B4-09-DS
% Mức: VD
Một vật dao động điều hòa trên trục Ox với hình vẽ
\choiceTF

\loigiai{
A. Sai — Tại $t = 0$, $x = 5\cos\!\left(-\dfrac{\pi}{3}\right) = 2,5$ cm, không phải $-2,5$ cm.
B. Sai — Tại $t = 1$ s, $x = 5\cos\!\left(\pi \cdot 1 - \dfrac{\pi}{3}\right) = 5\cos\!\left(\dfrac{2\pi}{3}\right) = -2,5$ cm, không phải $2,5$ cm.
C. Sai — Từ $t = 0$ đến $t = 1,5$ s, pha tăng từ $-\dfrac{\pi}{3}$ đến $\dfrac{7\pi}{6}$, vật đổi chiều chuyển động 2 lần, không phải 4 lần.
D. Đúng — Tại $t = \dfrac{11}{6}$ s, $x = 5\cos\!\left(\pi \cdot \dfrac{11}{6} - \dfrac{\pi}{3}\right) = 5\cos\!\left(\dfrac{3\pi}{2}\right) = 0$, nhưng đây là lần li độ bằng 0 thứ hai, không phải lần thứ nhất.
}
\end{ex}

% ===== Câu 34 =====
% ID: L11C1B4-09-TLN
% Mức: TH
\dangbt{Viết phương trình dao động điều hòa.}
\begin{ex}
% ID: L11C1B4-09-TLN
% Mức: TH
Vật dao động điều hòa có biên độ $4\,\text{cm}$ và quỹ đạo dài bao nhiêu cm?
\shortans{7/1/1900}
\end{ex}

% ===== Câu 38 =====
% ID: L11C1B4-10-TLN
% Mức: VD
\begin{ex}
% ID: L11C1B4-10-TLN
% Mức: VD
Vật dao động điều hòa có biên độ $6\,\text{cm}$ và quỹ đạo dài bao nhiêu cm?
\shortans{12}
\end{ex}

% ===== Câu 42 =====
% ID: L11C1B4-11-TLN
% Mức: VD
\begin{ex}
% ID: L11C1B4-11-TLN
% Mức: VD
Vật dao động điều hòa có biên độ $8\,\text{cm}$ và quỹ đạo dài bao nhiêu cm?
\shortans{15/01/1900}
\end{ex}

% ===== Câu 46 =====
% ID: L11C1B4-12-TLN
% Mức: TH
\dangbt{Xác định các đại lượng dựa vào đồ thị dao động ($x-t$, $v-t$, $a-t$).}
\begin{ex}
% ID: L11C1B4-12-TLN
% Mức: TH
Trên đồ thị li độ–thời gian, hai đỉnh dương liên tiếp ở $t=1$ s và $t=2$ s. Chu kì theo đơn vị s bằng bao nhiêu?
\shortans{31/12/1899}
\end{ex}

% ===== Câu 48 =====
% ID: L11C1B4-13-TLN
% Mức: VD
\begin{ex}
% ID: L11C1B4-13-TLN
% Mức: VD
Trên đồ thị li độ–thời gian, hai đỉnh dương liên tiếp ở $t=1$ s và $t=3$ s. Chu kì theo đơn vị s bằng bao nhiêu?
\shortans{2}
\end{ex}

% ===== Câu 50 =====
% ID: L11C1B4-14-TLN
% Mức: VD
\begin{ex}
% ID: L11C1B4-14-TLN
% Mức: VD
Trên đồ thị li độ–thời gian, hai đỉnh dương liên tiếp ở $t=1$ s và $t=5$ s. Chu kì theo đơn vị s bằng bao nhiêu?
\shortans{4}
\end{ex}

% ===== Câu 53 =====
% ID: L11C1B4-15-TN
% Mức: TH
\dangbt{Viết phương trình dao động điều hòa.}
\begin{ex}
% ID: L11C1B4-15-TN
% Mức: TH
Vật có biên độ $9\,\text{cm}$, chu kì $1\,\text{s}$ và ở biên dương lúc $t=0$. Phương trình đúng là
\choice
{\True $x=9\cos(2\pi t/1)$ cm}
{$x=9\sin(2\pi t/1)$ cm}
{$x=9\cos(\pi t/1)$ cm}
{$x=9\cos(2\pi t/1+\pi/2)$ cm}
\end{ex}

% ===== Câu 55 =====
% ID: L11C1B4-16-TN
% Mức: TH
\begin{ex}
% ID: L11C1B4-16-TN
% Mức: TH
Một vật dao động điều hoà dọc theo trục $Ox$ với phương trình $x=A\cos \omega t$. Nếu chọn gốc toạ độ $O$ tại vị trí cân bằng của vật thì gốc thời gian $t=0$ là lúc vật
\choice
{\True ở vị trí li độ cực đại thuộc phần dương của trục Ox}
{qua vị trí cân bằng $O$ ngược chiều dương của trục Ox}
{ở vị trí li độ cực đại thuộc phần âm của trục Ox}
{qua vị trí cân bằng $O$ theo chiều dương của trục Ox}
\loigiai{
Gốc thời gian $t=0$ là lúc vật ở vị trí li độ cực đại thuộc phần dương của trục Ox
}
\end{ex}

% ===== Câu 57 =====
% ID: L11C1B4-17-TN
% Mức: VD
\begin{ex}
% ID: L11C1B4-17-TN
% Mức: VD
Vật có biên độ $10\,\text{cm}$, chu kì $2\,\text{s}$ và ở biên dương lúc $t=0$. Phương trình đúng là
\choice
{\True $x=10\cos(2\pi t/2)$ cm}
{$x=10\sin(2\pi t/2)$ cm}
{$x=10\cos(\pi t/2)$ cm}
{$x=10\cos(2\pi t/2+\pi/2)$ cm}
\end{ex}

% ===== Câu 59 =====
% ID: L11C1B4-18-TN
% Mức: TH
\dangbt{Xác định các đại lượng dựa vào đồ thị dao động ($x-t$, $v-t$, $a-t$).}
\begin{ex}
% ID: L11C1B4-18-TN
% Mức: TH
Từ đồ thị dao động, khoảng thời gian giữa hai đỉnh dương liên tiếp là $1\,\text{s}$ và li độ cực đại là $4\,\text{cm}$. Chu kì bằng
\choice
{\True $1\,\text{s}$}
{$0.5\,\text{s}$}
{$2\,\text{s}$}
{$4\,\text{s}$}
\end{ex}

% ===== Câu 61 =====
% ID: L11C1B4-19-TN
% Mức: VD
\begin{ex}
% ID: L11C1B4-19-TN
% Mức: VD
Từ đồ thị dao động, khoảng thời gian giữa hai đỉnh dương liên tiếp là $2\,\text{s}$ và li độ cực đại là $6\,\text{cm}$. Chu kì bằng
\choice
{\True $2\,\text{s}$}
{$1\,\text{s}$}
{$4\,\text{s}$}
{$8\,\text{s}$}
\end{ex}

% ===== Câu 63 =====
% ID: L11C1B4-20-TN
% Mức: VD
\begin{ex}
% ID: L11C1B4-20-TN
% Mức: VD
Từ đồ thị dao động, khoảng thời gian giữa hai đỉnh dương liên tiếp là $4\,\text{s}$ và li độ cực đại là $8\,\text{cm}$. Chu kì bằng
\choice
{\True $4\,\text{s}$}
{$2\,\text{s}$}
{$8\,\text{s}$}
{$16\,\text{s}$}
\end{ex}
